\section{From $\ListLearn$ to the  $\kolp$ Problem}
\label{sec:kolp}

In this section, we show that for well-separated polytopes, a solution for the $\ListLearn$ problem can be used to solve the $\kolp$ problem as well. This algorithm uses the notion of soft convex hulls. 
We first describe the algorithm for constructing soft convex hulls, and then use it to solve the $\kolp$ problem. 

\subsection{Soft Convex Hulls}
\label{sec:soft}
Let $W$ be a finite set of points in $\Re^d$, and
$T$ be the vertices of $\ch(W)$. The subset $T$ of $W$ is the unique subset of $W$ with the following properties: 
\begin{itemize}
\item[(P1)] $W\subseteq \ch(T)$
\item[(P2)] $\forall w\in W$, if $w\notin \ch(W\setminus \{w\})$, then $w \in T$. 
\end{itemize}



We now define a natural notion of {\em soft} convex hull. 
\begin{defn}\label{soft-CH}
For an $\varepsilon \geq 0$, and $S\subseteq W$, define the $\varepsilon$-convex hull of $S$, $\ech(S)$, as $\ch(S) + \varepsilon \dia(W) B$, where $B$ is the unit ball of the Eucliean norm. 
\end{defn}

 The intuition behind the above definition is that $\ch(W)$ can have a very large set of vertices, but there may be a small set of points whose soft convex hull contains $W$. This is defined more formally as follows: 

\begin{defn}\label{envelope}
We call a subset $T\subseteq W$ an $\varepsilon$-envelope of $W$, written
$\eenv(W)$, if 
$W\subseteq \ech(T).$  
\end{defn} 

\noindent
{\bf Remarks:} The following observations about the set $\eenv(W)$ are easy to see:
\begin{itemize}
\item[(a)] There are several distinct sets $T$ which could qualify as $\eenv(W)$. For example, let $W$ consist of the following set of points in $\Re^2$:  a set of points $W_1$ close to $(0,0)$ and a set of points $W_2$ close to $(1,0)$. Let $T$ be a pair of points $\{x,y\}$ with $x \in W_1, y \in W_2$. Then it is easy to check that $T$ is an $\eenv(W)$. 
\item[(b)] Let $T$ be $\eenv(W)$. Unlike property~(P2) above, it is not necessary that if $w \in W$ is such that $w \notin \ech(W \setminus \{w\})$, then $w \in T$. 
 \end{itemize}

Since the set $\eenv(W)$ is not uniquely determined, we will impose one more condition on it to make it unique (if it exists) and polynomial time computable. This condition requires
the points of $T$ to be ``far apart'' from each other. More precisely:

\begin{defn}\label{strict-envelope}
For $\varepsilon,\delta \in [0,1]$, a set $T$ is called a $\edenv(W)$ if it is an $\eenv(W)$ and  
\begin{equation}\label{792}
\forall w\in T, \dist(w,\ch(T\setminus \{w\}))> \delta \dia(W)
\end{equation}
\end{defn}

\begin{fact}
\label{f:check}
Given a subset $T$ of $W$, we can check in polynomial time whether $T$ is an $\edenv(W)$. 
\end{fact}

\begin{proof}
Fix a subset $T$. We can verify in polynomial time whether $T$ is an $\eenv(W)$. Indeed, for each point $w \in W$, we need to check if $w \in \ech(T)$. This can be expressed as a convex program feasibility problem, where the variables are $\lambda_t$ for each $t \in T$: 
$$ \left|\sum_{t \in T} \lambda_t \cdot t - w \right| \leq \varepsilon \dia(W), \quad \sum_t \lambda_t = 1, \quad \lambda_t \geq 0 \  \  \forall t \in T.$$ 

Similarly, we can check~\eqref{792} using convex programming. For each $w \in T$, we can find $\dst(w, \ch(T \setminus  \{w\})$ as follows, where the variables are $\lambda_x, x \in T \setminus \{w\}$: 
$$ \min. \ \left| \sum_{x \in T \setminus \{w\}} \lambda_x \cdot x - w \right|, \quad \sum_x \lambda_x = 1, \quad \lambda_x \geq 0 \  \  \forall x \in T\setminus \{w\}. $$
\end{proof}


For rest of the section, we address the following question: given the set $W$, and parameters $\varepsilon, \delta$, is there a  $\edenv(W)$, and if so, can we find in polynomial time an approximation to this set? We informally argue that several ``natural'' greedy strategies do not work. Consider, for example, the following algorithms for identifying $\edenv(W)$: 
\begin{itemize}
    \item Identify $T$ as the set of points $w\in W$ for which $w$ is not close to $\ch(W\setminus \{w\})$: We define a set of points in $\Re^2$. We have {\em rings} of points, one around $ (0,0)$ and the other around $(0,1)$, each point is close to the convex hull of the others and so this algorithm will   dismiss all the points as not belonging to the desired set $T$.
    % \item \alert{This example is not clear to me, i.e., I am not sure all the conditions on the 4 points can be satisfied. - rk   }Start by adding an extreme point (along a randomly chosen direction) to $T$. For the next $k-1$ steps, iteratively find a point $w$ for which $\dist(w,T)$ is maximized and add it to $T$: We consider a set of points in $\Re^2$ as follows: let $v_1, v_2, v_3$ be three affinely independent points with $|v_1 - v_2| > |v_1 - v_3|/\delta, |v_2 - v_3|/\delta$ and $\dist(v_3, \ch(\{v_1, v_2\})) > \delta |v_1 - v_2|$. Now let $v_4$ be the mid-point of the line segment joining $(v_1, v_2)$. Let $W=\{v_1, v_2, v_3, v_4\}$. Say the first point selected by the greedy algorithm is $v_1$. Then the next two points selected by it will be $v_2, v_4$ respectively. But the ``correct'' $T$ is $\{v_1, v_2, v_3\}$. 
    \item Start by adding an extreme point (along an arbitrarily chosen direction) to $T$. For the next $k-1$ steps, iteratively find a point $w$ for which $\dist(w,\ch(T))$ is maximized and add it to $T$: a simple example as above shows that this idea may not work as well. Consider $v_1, v_2, v_3, v_4$ lying on corners of a square, and let $v_5$ be the mid-point of $v_3, v_4$. We start by adding $v_1$ and then $v_2$ to $T$. But in the next step, we could add $v_5$, and then both $v_3, v_4$. But the ``correct'' $T$ would have been $\{v_1, v_2, v_3, v_4\}$. 
\end{itemize}
 


If $\varepsilon=0$, the above question is easy to answer in polynomial time. The answer is yes iff the set $T$ of vertices of $\ch(W)$ satisfies (\ref{792}). Also, if $\delta=1$, $T$ has to be a singleton to satisfy (\ref{792}). 

In rest of this section, we consider the following problenm: For what pairs of values of $\varepsilon,\delta,$ can we prove that there is {\em essentially} at most one $\edenv(W)$, and if so, can we determine this set efficiently ? 
We do not know the exact answer to this, but our main result here (which suffices for the applications) is 
(verbally stated) an affirmative answer to the question if the following condition is satisfied:
$$\delta\in \Omega(\sqrt\varepsilon).$$
This will follow as a corollary of our main result:

\begin{theorem}
\label{thm:soft}
Let
$\delta,\varepsilon,\varepsilon_3$ be reals in $(0,1/8)$ satisfying
\begin{equation}
\label{822}\delta> \max \left(\frac{2\varepsilon}{\varepsilon_3-\varepsilon}, 4 \varepsilon_3 \right)
\end{equation}
Let $W$ be a finite set of points in $\Re^d$. 
We can 
determine in polynomial time whether  exists a set $T$ in $\edenv(W)$, and if so,  we can efficiently 
find a subset $Q$ of $W$ such that
\begin{align}
|Q|=|T|\\
\forall w\in T, \exists x\in Q: |w-x|\leq 2\varepsilon_3 \dia(W)
\end{align}
\end{theorem}
\begin{corollary}\label{cor:soft}
Let
$\delta,\varepsilon$ be reals in $(0,1/8)$ satisfying $\delta>16\sqrt\varepsilon$
Let $W$ be a finite set of points in $\Re^d$. 
We can 
determine in polynomial time whether  exists a set $T$ forming a $\edenv(W)$, and if so,  we can efficiently 
find a subset $Q$ of $W$ such that
\begin{align}
|Q|=|T|\\
\forall w\in T, \exists x\in Q: |w-x|\leq 8\sqrt \varepsilon \dia(W)
\end{align} 
\end{corollary}
\begin{proof}
The Corollary follows from Theorem (\ref{thm:soft}) by taking $\varepsilon_3=4\sqrt\varepsilon$ 
\end{proof}
\begin{proof} (of \Cref{thm:soft})
The procedure is described  in~\Cref{algo:main}. We first compute a subset $Q''$ of $W$ consisting of points $w \in W$ which do not lie in the soft convex hull of the points in $W$ which are ``far'' from $w$ -- this can be doen in polynomial time by using arguments similar to those in the proof of~\Cref{f:check}. Then $Q$ is defined as a maxumal subset of points in $Q''$ with the property that the pair-wise distance between the points in it large. Finally, we check whether $Q$ is an $\edenv(W)$, which can be done efficiently using~\Cref{f:check}. 


\begin{algorithm}[H]
  \caption{Procedure to identify $\edenv(W)$}
  \label{algo:main}
    Compute $\dia(W)$. \;
    Let $Q' := \{w \in W: w \in \ech \left(\{x \in W: |w-x| \geq \varepsilon_3 \dia(W) \}\right). $ \;
    Let $Q'' := W \setminus Q'$. \;
    Define $Q :=$ maximal subset of $Q''$ such that for every distinct $x,y \in Q$, $|x-y| > 2\varepsilon_3 \dia(W). $ \;
    \If{$Q$ is $\edenv(W)$} {
       Output  the set $Q$. \; }
    \Else {
     Output No. }
\end{algorithm}


Now we analyze this algorithm. If there is no $\edenv(W)$, then the algorithm will clearly say ``No'' (because the set $Q$ will not be $\edenv(W)$). So assume that there is a set $T$ which is $\edenv(W)$. 

\begin{claim}
\label{cl:ed}
$T \subseteq Q''$. 
\end{claim}
\begin{proof}

Suppose for the sake of contradiction that there is a point $w \in T \setminus Q''$.  For sake of brevity, let $W_x$ denote the subset of points in $W$ satisfying $|w-x|  \geq \varepsilon_3 \dia(W)$. The fact that $w \notin Q''$ implies that $w \in \ch(W_x)$, i.e., 
$$ w = \sum_{x \in W_x} \lambda_x \cdot x + e_0, $$
where $\lambda_x$ form a convex combination and $|e_0| \leq \varepsilon \dia(W).$ Since $W \subseteq \eenv(T)$, for each $x \in W_x$, there are a point $x' \in \ch(T)$ such that $e_x := x-x'$ has length at most $\varepsilon \dia(W)$. Since $w \in T$, we can write $x'$ as 
$$ x' = \mu_x w + (1-\mu_x) y_x, \quad y_x \in \ch(T \setminus \{w\}), $$
where $\mu_x \in [0,1]$. 
Therefore, 
$$ |w-x| \leq |w-x'| + |e_x| = (1-\mu_x) |w-y_x| + |e_x| \leq (1-\mu_x) \dia(W) + \varepsilon \dia(W). $$
But we know that $|w-x| \geq \varepsilon_3 \dia(W)$. Therefore, we get
\begin{align}
\label{eq:111}
(1- \mu_x) \geq \varepsilon - \varepsilon_3. 
\end{align}
Now, 
\begin{align*}
w & = \sum_{x \in W_x} \lambda_x \cdot x + e_0 = \sum_{x \in W_x} \lambda_x \cdot x' + \sum_x \lambda_x e_x + e_0 \\
& = \sum_{x \in W_x} \lambda_x \mu_x \cdot w + \sum_{x \in W_x} \lambda_x (1 - \mu_x) y_x + e,  
\end{align*}
where $e = \sum_x \lambda_x e_x + e_0$ has length at most $2\varepsilon \dia(W)$. 
Let $\theta_x$ denote $\lambda_x(1-\mu_x)$. Observe that $\theta:= \sum_x \theta_x = 1 - \sum_x \lambda_x \mu_x$. Therefore, we can rewrite the above as 
$$ w - \frac{1}{\theta} \sum_{x \in W_x} \theta_x y_x = \frac{e}{\theta}. $$
Since $\frac{1}{\theta} \sum_{x \in W_x} \theta_x y_x \in \ch(T \setminus \{w\})$, we see that 
$$ \dist(w, \ch(T \setminus \{w\}) \leq \frac{|e|}{\theta} \leq 
\frac{2\varepsilon \dia(W)}{\sum_{x \in W_x} \lambda_x (1-\mu_x)} 
\stackrel{\eqref{eq:111}}{\leq} \frac{2\varepsilon}{\varepsilon-\varepsilon_3} \dia(W). $$
Using~\eqref{792} and~\eqref{822}, we get a contradiction. Therefore, $w \in Q''$. 
\end{proof}

Since $T \subseteq Q''$, for every $w \in T$, there exists a point $x_w \in Q$ such that $|w-x_w| \leq 2 \varepsilon_3 \dia(W). $ It is easy to see that if $w, w'$ are two distinct points in $T$, then $x_w \neq x_{w'}$. Indeed, 
$$ |x_w - x_{w'}| \geq |w-w'| - |w-x_w| - |w'-x_{w'}| \geq \delta \dia(W) - 4 \varepsilon_3 \dia(W) \stackrel{\eqref{822}}{>} 0. $$
Finally, we prove that every element in $Q$ is of the form $x_w$ for some $w \in T$. 

\begin{claim}
\label{cl:final}
Every $x \in Q$ is of the form $x_w$ for some $w \in T$. 
\end{claim}
\begin{proof}
Consider a point $x \in Q$. This implies that $x \in Q''$. We claim that there is a point $w \in T$ such that $|x-w| \leq \varepsilon_3 \dia(W)$. Suppose not. Then the set of points $w$ such that $|x-w| \geq \varepsilon_3 \dia(W)$ includes $T$. But we know that $x \in \ech(T)$, and so $x \in Q'$, a contradiction. 

Therefore $$|x - w_x| \leq |x- w| + |w-w_x| \leq  2\varepsilon \dia(W). $$
This implies that $x=w_x$ (since any two distinct points in $Q$ have distance greater than $2\varepsilon \dia(W)$. 
\end{proof}

Thus we have shown that $|Q|=|T|$ and for every $w \in T$, there is a unique element $x_w \in Q$ with $|w-x_w| \leq 2 \varepsilon_3 \dia(W)$.
\end{proof}




\subsection{Algorithm for $\kolp$}
We now show how soft convex hulls can be used to generate a solution for the $\kolp$ problem. The following result, which formalizes~\Cref{kolp-0}, uses the same setting as that in~\Cref{upper-bound-1}:

\begin{theorem}\label{upper-bound-kolp}
Suppose $\varepsilon,\delta$ are reals in $[0,1]$ with $\delta^2\geq c\varepsilon\sqrt d, \delta^3 \geq c \, \varepsilon,$ where $c$ is a large enough constant.  Let $K$ be a $\delta$-well-separated $k$-vertex polytope. Suppose we are also given an  $ \opto_\varepsilon(K)$ oracle $\cal O$. 
Let $W$ be the set of answers of the oracle $\cal O$ to $m=poly(d)\cdot k^{\Omega(1/\delta^2)}$ independent random queries. We can find $Q\subseteq W, |Q|=k$ in 
randomized $poly(d) \cdot k^{\Omega(1/\delta^2)}$-time which satisfies the following condition w.h.p.:  for every vertex $v$ of $K$, there is a point $v'$ in $Q$ with $|v-v'| \leq \delta \Delta(K)/10$. 
\end{theorem}
\begin{proof}
    The proof of~\Cref{upper-bound-1} shows that  for every vertex  $\Mdot{\ell}$ of $K$, there is point $x(u^{\ell}) \in W$ such that 
    \begin{align}
        \label{eq:distMl}
          |x(u^{(\ell)})-M_{\cdot,\ell}|\leq 17 \delta^2 \Delta(K)/c.
    \end{align}
    Let $T$ denote the set of points $\{x(\uell): \ell \in [k]\}$. Our first claim is that the points of $T$ are also well-separated. 

\begin{claim}
\label{cl:wellstx}
For any $\ell \in [k]$, 
$$\dst(x(\uell), \CH(T \setminus \{x(\uell)\})) \geq \delta \Delta(K)/2.$$ Further, the diameter of $\CH(T)$ is at most $2 \Delta(K). $
\end{claim}
\begin{proof}
    Fix an index $\ell \in [k]$ and a point $y \in \CH(T \setminus \{x(\uell)\})).$ We can express $y$ as a convex  combination of points in $T \setminus \{x(\uell)\},$ i.e., 
    $$ y = \sum_{\ell' \neq \ell} \lambda_{\ell'} \cdot x(u^{(\ell')}), \quad \mbox{where \,  } \sum_{\ell' \neq \ell} \lambda_{\ell'} = 1.$$ Now
    \begin{align*}
        |x(\uell) - y| & \geq \left| \Mdot{\ell} - \sum_{\ell' \neq \ell} \lambda_{\ell'} \cdot \Mdot{\ell'} \right| - |x(\uell) - \Mdot{\ell}| - \sum_{\ell' \neq \ell} \lambda_{\ell'} \cdot |x(u^{(\ell')})-\Mdot{\ell'}| \\
        & \geq \left(\delta - \frac{34 \delta^2}{c} \right) \Delta(K) \geq \delta \Delta(K)/2,
    \end{align*}  
    where the second last inequality follows from~\eqref{eq:distMl} and the fact that $K$ is $\delta$-well-separated. 
    Since $\dst(x(\uell), K) \leq \varepsilon \Delta(K)$, it follows that $\Delta(\CH(T)) \leq 2 \Delta(K).$
    %A similar argument shows that for any two distinct $\ell, \ell' \in [k]$, 
    %$$ |x(\uell) - x(u^{(\ell')})| \leq (1+ \delta)\Delta(K) %\leq 2 \Delta(K).$$
\end{proof}

Recall that $W$ denotes the set $\{x(u): u \in U\}$. We now show that $\CH(T)$ closely approximates the set $\CH(W)$. 
\begin{claim}
    \label{cl:T}
    $W \subseteq \echp(T), $ where $\varepsilon'= \frac{32 \delta^2}{c}. $
\end{claim}

\begin{proof}
    Fix a point $x(u) \in W$. We know that $x(u)$ can be written as 
    $$ x(u) = y(u) + z(u), \quad y(u) \in K, |z(u)| \leq \varepsilon.$$
    Let $y(u) = \sum_{\ell \in [k]} \lambda_{\ell} \cdot \Mdot{\ell}, $ where the coefficients $\lambda_\ell$ form a convex combination. Then 
    $$ \left|x(u) - \sum_{\ell \in [k]} \lambda_\ell \cdot x(\uell) \right| \leq \sum_{\ell \in [k]} \lambda_\ell \cdot |x(\uell) - \Mdot{\ell}| + |z(u)| \leq \frac{17 \delta^2 \Delta(K)}{c} + \varepsilon \Delta(K) \leq \frac{32 \delta^2 \Delta(K)}{c},  $$
    where the second last inequality follows from~\eqref{eq:distMl} and the last inequality by the assumption in. This proves the desired result. 
\end{proof}


~\Cref{cl:wellstx} and~\Cref{cl:T} imply that $T$ is $\edenvp(W)$ with $\delta' = \delta/4, \varepsilon'= \frac{32 \delta^2}{c}$. We can now apply~\Cref{cor:soft} to get approximations to $x(\uell)$ within distance $17 \sqrt{\varepsilon'} \Delta(K)$.~\Cref{cl:T} now implies that we can get approximations to $\Mdot{\ell}$ within distance 
$$ 17 \sqrt{\varepsilon'} \Delta(K) + \frac{16 \delta^2 \Delta(K)}{c} \leq \frac{\delta \Delta(K)}{10}.$$
This proves the desired result.
\end{proof}

The following result is the analogue of~\Cref{upper-bound-2}, and gives an algorithm for the $\kolp$ problem under slightly different conditions on the parameters $\delta$ and $\varepsilon$. 

\begin{theorem}\label{upper-bound-2kolp}
Suppose $\varepsilon,\delta$ are reals in $[0,1]$ with $\delta^2\geq c\varepsilon\sqrt d, \delta \geq \frac{ \sqrt{\log k}}{\sqrt{cd}},$ where $c$ is  large enough constant.  Let $K$ be a $\delta$-well-separated $k$-vertex polytope. Suppose we are also given an $\opto_\varepsilon(K)$ oracle $\cal O$.
Let $W$ be the set of answers of the oracle $\cal O$ to $m=poly(d)\cdot k^{O(1/\delta^2)}$ independent queries.  In randomized $poly(d) \cdot k^{O(1/\delta^2)}$-time we can find $Q\subset W$ of $k$ points such that the following condition is satisfied w.h.p.:  for every vertex $v$ of $K$, there is a point $v'$ in $Q$ with $|v-v'| \leq \delta \Delta(K)/10$. 
\end{theorem}

% \subsection{Verbal Description of results}

% We formulate a problem on convex sets described by an oracle. The study of convex sets described only by oracles which respond to queries about the set is an important part of Convex Geometry, starting with the seminal work of~\cite{GLS1988}. In particular, Membership oracles, Separation oracles and Optimization oracles are well-studied and in summary, the central results of~\cite{GLS1988} and following literature establish polynomial time reductions among these oracles under mild assumptions. 

% Here, we consider a different oracle 
% from the ones which have been studied: Our convex set is a (possibly low dimensional) unknown non-empty polytope $K$ with $k$ vertices in ${\bf R}^d$ and on a queried unit length vector $u\in {\bf R}^d$, the oracle returns an approximately feasible and approximately optimal solution to the linear program: $$\mbox{Maximize } u\cdot x\; ;\; x\in K.$$ The problem we wish to solve, called the {\color{red} Oracle Latent $k$-Polytope ($\kolp$)} problem is to find approximations to the vertices of $K$ using of oracle calls and other computations bounded above by a polynomial in $d,k$ and error parameters. [More formal definitions of these quantities and $\kolp$ are given below.]
% Unlike many well-studied oracles admitting polynomial time reductions, we will see  
% %
% \begin{itemize}
%     \item [a] via an Information-theoretic argument that if the approximation error (defined formally below) is $\Omega^*(1/\sqrt d)$, a polynomial time algorithm is impossible and further, 
%     \item [b] making crucial use of our new Random Separating Hyperplane theorem, we will give a polynomial time algorithm when the error is $O(1/\sqrt d)$.
% \end{itemize}
% Thus, this provides a nearly-tight bound on error required for poly time computation with the model. 

% ??????????BELOW BIT TO BE PROPERLY EDITED??????????

% But this is not the end of the story because the best-known polynomial time computable oracles for LkP have errors which are $O(1/k^c)$ and if $k^c<<\sqrt d$,
% then, the error is too high to use b.

% ???CHECK - SPELL OUT $c$????

% We will overcome this making a new use of a well-used method - Projecting to the SVD subspace. In the original version of LkP (not the oracle version here -??? better wording  ???),  ??? given $\bA$ ?? latent $\bP$  we can show that if we project data to the $k-$SVD subspace, then, $\hat K$ the projection has good vertex separation and also we have a $O(1/sqrt k)$ oracle....????  

% ???????????END OF TO BE EDITED???????

% \subsection{Formal Definition of $\kolp$}

% The underlying problem of finding the vertices of a latent polytope was formulated in SODA ??? and called LkP. The oracle formulation here is an abstraction of LkP. LkP as well as our OLP are interesting objects of study because 
% they include as special cases a number of problems from Machine Learning - among them, Learning Stochastic Mixtures, Clustering, Stochastic Block Models, Topic Models and Mixed Membership Community Models, as we describe later in the paper. 




% Now we define the $\kolp$ problem formally. 



% %We henceforth assume that $\Delta(K)$ is at most 1. Note that there is no loss of generality since this can be ensured by scaling. 
% {\color{red} OMIT THE FOLLOWING DEFINITION FROM HERE> ALREADY IN THE INTRO}
% \begin{defn}\label{OLP}
% Let $K$ be a  $\delta$-well-separated polytope in $\Re^d$ with $k$ vertices $\Mdot{1}, \Mdot{2}, \ldots, \Mdot{k}$.  An {\em $\varepsilon$-optimization} oracle for $K$ takes as input  a unit length vector $u\in \Re^d$ and returns a vector $x\in \Re^d$ satisfying 
% $$x\in K+\varepsilon \Delta(K) B_d;\quad u\cdot x\geq \mbox{Max}_{y\in K} \, u\cdot y-\varepsilon \Delta(K),$$ 
% where $B_d$ denotes the unit ball in $\Re^d$. 

% The {\em $\kolp$ problem} is to 
% find a set $T$ of $k$ points in $\Re^d$ approximating the vertices of $K$, i.e., 
% $$\forall \ell\in [k], \, \exists x\in T, \, 
% |x-M_{\cdot,\ell}|\leq \delta \Delta(K)/4.$$
% \end{defn}

% We assume $\delta> 2\varepsilon$. This implies that any $T$ as above must have a distinct point approximating each vertex of $K$.
% We will presently see that if $$\varepsilon \in\Omega\left( \frac{ \ln d}{\sqrt d}\right),$$ then no polynomial time algorithm can solve the $\kolp$ problem
% (see Theorem (\ref{lower-bound})). 
% We prove this below with what we call the {\em needle in a $d$-dimensional haystack} example.

% In contrast we will prove an upper bound which is tight within logarithmic factors: namely, that if we had an $O(1/\sqrt d)$-optimization oracle, then, indeed, a collection of randomly chosen directions can be used to approximately identify each vertex of $K$ (See Theorem (\ref{upper-bound-1}). The proof of Theorem (\ref{upper-bound-1}) crucially uses the Random Separating Hyperplane Theorem (\Cref{thm:sephp}). 

% Consider a simpler problem than $k-OLP$ (which we denote ``Simple-$k-OLP$) where, one is required to output a point which is within $\varepsilon \Delta (K)$ of some vertex of $K$. We prove:

% \begin{theorem}\label{lower-bound-thm}
%     The Simple-$k-OLP$ problem with $\varepsilon\geq 8\ln d/\sqrt d$ cannot be solved in deterministic polynomial time even when $\delta =1/2$.
% \end{theorem}

% \begin{proof}
%     The proof is by producing an ``adversarial oracle''. Our $K$ will a ``needle'' of the form $\{\lambda u:\lambda\in [-1,1]\}$, where, $u$ is a unit length vector. The oracle's answer for each query will be always the 0 vector. Suppose $v_1,v_2,\ldots, v_q, q=d^c$ for a constant $c$ be the vectors on which the oracle is queried. The answer 0 is clearly consistent with $u$ being any vector with $|u\cdot v_i|\leq 4\ln d/\sqrt d$. For a single $v_i$, the probability that a random $u$ satisfies $|u\cdot v_i|>4\ln d/\sqrt d$ is at most $e^{-d\varepsilon^2}$. So by union bound, we see that there is a vector $u_1$ satisfying 
%     $$|u_1\cdot v_i|\leq 4\ln d/\sqrt d\forall i\in [q].$$
%     Further, by a similar argument, there is a vector $u_2$satisfying
%     $$|u_2\cdot v_i|\leq 4\ln d/\sqrt d, i\in [q]\; ;\; |u_1-u_2|,|u_1+u_2|\geq .1.$$ It follows that the two possible $K$, where 
%     $$\{ \lambda u_1:\lambda\in [-1,1]\}\; ,\; \{ \lambda u_2:\lambda\in [-1,1]\}$$
% are both consistent with the answer 0 for all $v_i,i\in [q]$ and in addition, we have that no point $w\in {\bf R}^d$ is within distance $\varepsilon$ of two vertices, one from each of these 2 needles. So no answer is valid for both needles and the adversarial oracle can choose one of the needles to render the algorithm's answer incorrect.
% \end{proof}


% \begin{theorem}\label{lower-bound}
%          The $\kolp$ problem with $\varepsilon \geq  \frac{8 \ln d}{\sqrt d},$ cannot be solved in polynomial time even when $\delta=1/2$.
% \end{theorem}

% \begin{proof}
% We give the lower bound proof for a deterministic algorithm, though the lower bound can be easily extended to randomized algorithms. Let $\cal A$ be any deterministic algorithm for the $\kolp$ problem, and let the running time of $\cal A$ be at most $d^c$, for a constant $c$ (the input polytope will have at most $d$ vertices). The algorithm $\cal A$ has access to an $\varepsilon$-optimization oracle $\cal O$ for a latent polytope $K$. In our lower bound construction below, $K$ will have exactly 2 vertices, and we shall show that with high probability, $\cal A$ will not be able to output a point within $\varepsilon \Delta(K)$ of $K$.  


% The proof is by producing an ``adversarial oracle''. Our $K$ will a ``needle'' of the form $\{\lambda u:\lambda\in [-1,1]\}$, where, $u$ is a unit length vector. The oracle's answer for each query will  always be the 0 vector. Suppose $v_1,v_2,\ldots, v_q,$ where $q=d^c$,  be the vectors on which the oracle is queried. The answer 0 is clearly consistent with $u$ being any vector with $|u\cdot v_i|\leq 4\ln d/\sqrt d$. For a single $v_i$, the probability that a random $u$ satisfies $|u\cdot v_i|>4\ln d/\sqrt d$ is at most $e^{-d\varepsilon^2}$. So by union bound, we see that there is a vector $u_1$ satisfying 
%     $$|u_1\cdot v_i|\leq 4\ln d/\sqrt d \quad \quad \forall i\in [q].$$
%     Further, by a similar argument, there is another unit vector $u_2$ satisfying
%     $$ |u_1-u_2|,|u_1+u_2|\geq 0.1, \quad \mbox{and} \quad |u_2\cdot v_i|\leq 4\ln d/\sqrt d \quad  \forall i\in [q].$$ It follows that both the following  two choices for $K$
%     $$\{ \lambda u_1:\lambda\in [-1,1]\}\; ,\; \{ \lambda u_2:\lambda\in [-1,1]\}$$
% are consistent with the answer 0 for all $v_i,i\in [q],$ and in addition, we have that no point $w\in {\bf R}^d$ is within distance $\varepsilon$ of two vertices, one from each of these two needles. So no answer is valid for both the needles and the adversarial oracle can choose one of the needles to render the algorithm's answer incorrect.

% The proof is based on the usual  advesarial oracle model ???EXPLAN ADVERSARIAL ORACLE??? .
% The set-up we use to prove the lower bound may be called the {\color{red} Needle in a $d-$dimensional haystack} problem: 
% We give the lower bound proof for a deterministic algorithm, though the lower bound can be easily extended to randomized algorithms. Let $\cal A$ be any deterministic algorithm for the $$ problem. The algorithm $\cal A$ has access to an $\varepsilon$-optimization oracle $\cal O$ for a latent polytope $K$. In our lower bound construction below, $K$ will be have at most 2 vertices. 

% Now, the algorithm $\cal A$ can make at most $q:=d^c$ queries to $\cal O$, where $c$ is a constant. Whenever $\cal A$ makes a query $v$, where $v$ is a unit vector, the oracle $\cal O$ outputs the vector 0. We now show that there are two distinct polytopes which satisfy the answers of the oracle $\cal O$. 

% The first polytope $K_1$ is simply the polytope consisting of the single point $\{0\}$ (since we want the parameter $k$ to be fixed and known, we could think of $K_1$ as a 1-dimensional polytope with two vertices very close to the 0 vector). We now describe the second polytope $K_2$. Let $v_1, \ldots, v_q$ be the vectors queried by the algorithm $\cal A$. We claim  that there is a unit vector $u$ such that for each $i \in [q]$, 
% \begin{align}
% \label{eq:advcond}
%      |u \cdot v_i| \leq \frac{4 \ln d}{\sqrt{d}}.
% \end{align}

% Indeed, for a fixed vector $v_i$, the  measure of the set of unit vectors not satisfying the above condition is at most $\exp(-d \varepsilon^2) < \frac{1}{d^c}$ for large enough $d$. A simple  union bound over the $q$ query vectors now implies that there exists a unit vector $u$ satisfying the above condition for each of the query vectors $v_i$. 

% We define the second polytope $K_2$ as the 1-dimensional polytope: $\{ \lambda u: -1 \leq \lambda \leq 1\}$. We first claim  that the answers of $\cal O$ are consistent with both $K_1$ and $K_2$. The fact that they are consistent with $K_1$ is easy to see. For $K_2$, condition~\eqref{eq:advcond} implies that each of the queries $v_i$ has optimal value at most $\frac{4 \ln d}{\sqrt{d}}.$
% Since $\varepsilon \geq \frac{8 \ln d}{\sqrt{d}}, $ the oracle's answer of 0 is consistent with the permitted error bounds. 

% Now suppose $\cal A$ outputs the vertices $w_1, w_2$. If either of $w_1, w_2$ is within distance $1/2$ of the $0$ vector, we set $K_2$ as the latent polytope, otherwise we set $K_1$ as the latent polytope.  In the second case, the distance between $\{w_1, w_2\}$ and the vertices of $K_1$ is at least $1/2$. In the first case, suppose $w_1$ is within distance $1/2$ of $0$. Then $w_1$ is at least distance $1/2$ away from the vertices of $K_2$. The vertex $w_2$ must also be at least $1/2$ away from one of the vertices of $K_2$. Thus, $\cal A$ cannot approximate both the vertices of $K_2$ with distance $1/2$. This proves the desired result. 










% The adversary has a Latent Polytope which is either $\{ 0\}$ (consisiting of just the single point - the origin) or a 1-d polytope (an interval) $\{ \lambda u:\lambda \in [-1/2\; ,\; 1/2]\},$ where, $u$ is an unknown unit length vector. 

% For every query $v$, the adversary  answers $x=\underline 0$. The answer rules out any $u$ with $|u\cdot v|\geq 4\ln d/\sqrt d$ (for such $u$ 's, 0 is not a valid answer, but for $u$ 's satisfyng $|u\cdot v|\leq 4\ln d/\sqrt d$, 0 is a valid answer.)

% If $s\leq $ poly many queries are asked, say the query vectors are $v_1,v_2,\ldots ,v_s$. All $u$ satisfying $|u\cdot v_t|\leq \varepsilon,$ for $t=1,2,\ldots ,s$ are consistent with the answer 0, so far. Now, the measure of $u$ with $|u\cdot v_t|>\varepsilon$ for one $t$ is at most
% $\exp(-d\varepsilon^2)$ and thus by union bound, we have at least one consistent $u$ left after poly many queries. Also the single point polytope $\{ 0\}$ is also consistent. So there is not unique answer 
% yet and the algorithm will fail.
%\end{proof}
% \section{Algorithm for $\kolp$ in $\Re^d$}
% \label{sec:kolp}

% We first present an algorithm which needs $\varepsilon\in O^*(1/\sqrt d)$. In the next section, we will see that we can project to the $k-$dimensional SVD subspace and solve the problem in the projection, thus solving it assuming only
% $$\varepsilon\in O^*(1/\sqrt k).$$
% %
% {\color{red} In the algorithm , we need to know $\delta$, but, $\varepsilon$ need not be. In general, should we write down the parameters any alg has to know?}
% %

%\begin{proof}
%  The algorithm chooses a set $U$ of $k^{\Omega(1/\delta^2)}$ unit length i.i.d. Gaussian vectors. For each $u\in U$, it calls the $\varepsilon$-optimization oracle to find a vector $x(u)$. Let $W$ denote the set $\{x(u): u \in U\}$. We invoke~\Cref{cor:soft} with the parameters: 
% $$ \delta' := \delta/4, \, \varepsilon' := \frac{32 \delta^2}{c}$$
% on the set $W$. The algorithm in that result outputs a set $Q$. We output that $Q$ as the approximation to the set of vertices of $K$. We now prove that the set $Q$ has the desired properties. We first show that for every vertex of $K$, there is a direction $u$ in $U$ along which the projection of this vertex is higher than the projection of the remaining vertices by a large enough margin. 
% Let $\Mdot{1}, \ldots, \Mdot{k}$ be the vertices of $K$. 
% \begin{claim}
%     \label{cl:ul}
%     With high probability, the following event happens: for each $\ell \in [k]$, there is a vector $u^{(\ell)} \in U$ such that for all $\ell' \in [k], \ell' \neq \ell,$ we have
%     \begin{align}
%         \label{eq:100}
%         u^{(\ell)}\cdot M_{\cdot, \ell}> u^{(\ell)}\cdot M_{\cdot, \ell'}+\frac{c \, \varepsilon \cdot \Delta(K)}{8 \delta^2}
%     \end{align}
% \end{claim}

% \begin{proof}
%     Fix a vertex $\Mdot{\ell}$ of $K$.  Let $K'$ be the convex hull of $\{\Mdot{1}, \ldots, \Mdot{k} \} \setminus \{\Mdot{\ell}\}.$ We invoke~\Cref{thm:sephp} on the polytope $K'$ and the point $a := \Mdot{\ell}$. The definition of $\delta$-well-separated implies that~\eqref{eq:condhp} is satisfied with $\rho = \delta$. Using $V = \Re^d$ in the statement of~\Cref{thm:sephp}, we see that $$\Pr_{u \in U} \left[ u \cdot \Mdot{\ell} - \max_{\ell' \in [k], \ell' \neq \ell} u \cdot \Mdot{\ell'} \geq \frac{\delta \sqrt{\log k} \Delta(K)}{\sqrt{\log k} + 4 \delta \sqrt{d}} \right] \geq \frac{1}{40} k^{-10/\delta^2}.$$
%     Since 
%     $$\frac{\delta \sqrt{\log k}}{\sqrt{\log k} + 4 \delta \sqrt{d}}  \geq \min \left( \frac{\delta}{2}, \frac{1}{8 \sqrt{d}} \right) \geq \frac{\varepsilon c}{8 \delta^2}, $$
%     the desired result follows from the fact that $|U| \gg \frac{1}{40} k^{-10/\delta^2}.$
    
% \end{proof}

% For rest of the proof, assume that the statement in~\Cref{cl:ul} holds true, i.e., there are directions $u^{(1)}, \ldots, u^{(k)} \in U$ satisfying~\eqref{eq:100}. 
% We now show that for every vertex of $\Mdot{\ell}$ of $K$, the corresponding point $x(u^{\ell})$ is close to  $\Mdot{\ell}$. 
% \begin{claim}
%     \label{cl:xl}
%     For every $\ell \in [k]$, 
%     $$|x(u^{(\ell)})-M_{\cdot,\ell}|\leq 17 \delta^2 \Delta(K)/c.$$
% \end{claim}
% \begin{proof}
%     By the definition of $\varepsilon$-optimization oracle, we know that $x(u^{(\ell)})$ can be written as $y(u^{\ell})+z(\uell)$, where $y(\uell) \in K$ and $|z(\uell)| \leq \varepsilon \Delta(K)$. Thus, there is a convex combination $\lambda_{\ell'}, \ell' \in [k],$ of the vertices $\Mdot{\ell'}$ of $K$ such that 
%     $$ x(\uell) = \sum_{\ell' \in [k]} \lambda_{\ell'} \Mdot{\ell'} + z(\uell).$$
   

% By the definition of the $\varepsilon$-optimization oracle, $x(u^{(\ell)})\cdot u^{(\ell)}\geq M_{\cdot, \ell}\cdot u^{(\ell)}-\varepsilon\Delta (K)$ and $|z(u^{(\ell)})\cdot u^{(\ell)}|\leq\varepsilon\Delta(K)$. So, we get
% $$M_{\cdot,\ell}\cdot u^{(\ell)}-\varepsilon\Delta(K)\leq \sum_{\ell' \in [k]}\lambda_{\ell'} M_{\cdot,\ell'}\cdot u^{(\ell')}+\varepsilon\Delta(K),$$
% which implies (after subtracting $\lambda_\ell M_{\cdot,\ell} \cdot u^{(\ell)}$ from both sides):
% $$(1-\lambda_\ell) M_{\cdot,\ell}u^{(\ell)}\leq \sum_{\ell'\not= \ell} \lambda_{\ell'} M_{\cdot,\ell'}\cdot u^{(\ell')}+2\varepsilon\Delta(K)$$
% which, using~\Cref{cl:ul}, yields:
% $$(1-\lambda_\ell) M_{\cdot,\ell}u^{(\ell)}\leq (1-\lambda_\ell) M_{\cdot,\ell}u^{(\ell)}-(1-\lambda_\ell)\frac{c\varepsilon\Delta(K)}{8\delta^2}+2\varepsilon\Delta(K).$$
% %which leads to:


    
% It follows from the above inequality that 
% \begin{align}
% \label{eq:105}
%     1- \lambda_\ell \leq \frac{16 \delta^2}{c}.
% \end{align}
%     Therefore, 
% \begin{align}
%     |x(\uell) - \Mdot{\ell}| = \left| \sum_{\ell' \neq \ell} \lambda_{\ell'} (\Mdot{\ell} - \Mdot{\ell'})  
%     \right| + \varepsilon  \Delta(K) \leq \sum_{\ell' \neq \ell} \lambda_{\ell'} \Delta(K) + \varepsilon \Delta(K) \stackrel{\eqref{eq:105}}{\leq} \frac{17 \delta^2 \Delta(K)}{c}.
%     \notag
% \end{align}
% %where the second inequality follows from the fact that $\Delta(K) \leq 1$. 
% \end{proof}

% Let $T$ denote the set of points $\{x(\uell): \ell \in [k]\}$. Our first claim is that the points of $T$ are also well-separated. 

% \begin{claim}
% \label{cl:wellstx}
% For any $\ell \in [k]$, 
% $$\dst(x(\uell), \CH(T \setminus \{x(\uell)\})) \geq \delta \Delta(K)/2.$$ Further, the diameter of $\CH(T)$ is at most $2 \Delta(K). $
% \end{claim}
% \begin{proof}
%     Fix an index $\ell \in [k]$ and a point $y \in \CH(T \setminus \{x(\uell)\})).$ We can express $y$ as a convex  combination of points in $T \setminus \{x(\uell)\},$ i.e., 
%     $$ y = \sum_{\ell' \neq \ell} \lambda_{\ell'} \cdot x(u^{(\ell')}), \quad \mbox{where \,  } \sum_{\ell' \neq \ell} \lambda_{\ell'} = 1.$$ Now
%     \begin{align*}
%         |x(\uell) - y| & \geq \left| \Mdot{\ell} - \sum_{\ell' \neq \ell} \lambda_{\ell'} \cdot \Mdot{\ell'} \right| - |x(\uell) - \Mdot{\ell}| - \sum_{\ell' \neq \ell} \lambda_{\ell'} \cdot |x(u^{(\ell')})-\Mdot{\ell'}| \\
%         & \geq \left(\delta - \frac{34 \delta^2}{c} \right) \Delta(K) \geq \delta \Delta(K)/2,
%     \end{align*}  
%     where the second last inequality follows from~\Cref{cl:xl} and the fact that $K$ is $\delta$-well-separated. 
%     Since $\dst(x(\uell), K) \leq \varepsilon \Delta(K)$, it follows that $\Delta(\CH(T)) \leq 2 \Delta(K).$
%     %A similar argument shows that for any two distinct $\ell, \ell' \in [k]$, 
%     %$$ |x(\uell) - x(u^{(\ell')})| \leq (1+ \delta)\Delta(K) %\leq 2 \Delta(K).$$
% \end{proof}

% Recall that $W$ denotes the set $\{x(u): u \in U\}$. We now show that $\CH(T)$ closely approximates the set $\CH(W)$. 
% \begin{claim}
%     \label{cl:T}
%     $W \subseteq \echp(T), $ where $\varepsilon'= \frac{32 \delta^2}{c}. $
% \end{claim}

% \begin{proof}
%     Fix a point $x(u) \in W$. We know that $x(u)$ can be written as 
%     $$ x(u) = y(u) + z(u), \quad y(u) \in K, |z(u)| \leq \varepsilon.$$
%     Let $y(u) = \sum_{\ell \in [k]} \lambda_{\ell} \cdot \Mdot{\ell}, $ where the coefficients $\lambda_\ell$ form a convex combination. Then 
%     $$ \left|x(u) - \sum_{\ell \in [k]} \lambda_\ell \cdot x(\uell) \right| \leq \sum_{\ell \in [k]} \lambda_\ell \cdot |x(\uell) - \Mdot{\ell}| + |z(u)| \leq \frac{17 \delta^2 \Delta(K)}{c} + \varepsilon \Delta(K) \leq \frac{32 \delta^2 \Delta(K)}{c},  $$
%     where the second last inequality follows from~\Cref{cl:xl} and the last inequality by the assumption in. This proves the desired result. 
% \end{proof}


% ~\Cref{cl:wellstx} and~\Cref{cl:T} imply that $T$ is $\edenvp(W)$ with $\delta' = \delta/4, \varepsilon'= \frac{32 \delta^2}{c}$. We can now apply~\Cref{cor:soft} to get approximations to $x(\uell)$ within distance $17 \sqrt{\varepsilon'} \Delta(K)$.~\Cref{cl:T} now implies that we can get approximations to $\Mdot{\ell}$ within distance 
% $$ 17 \sqrt{\varepsilon'} \Delta(K) + \frac{16 \delta^2 \Delta(K)}{c} \leq \frac{\delta \Delta(K)}{10}.$$
% This proves the desired result. 
% \end{proof}

% The statement of~\Cref{upper-bound-1} requires two different bounds relating $\delta$ to $\varepsilon$. In some applications, it may be difficult to ensure both of these conditions. We consider the setting when $d \gg k$, where the following variation of~\Cref{upper-bound-1} is better suited. 


% \begin{proof}
%     The proof is identical to that of~\Cref{upper-bound-1}, except that in the proof of~\Cref{cl:ul}, we now have the following modified argument. Since $\delta \sqrt{cd} \geq \sqrt{\log k}, $ we get 
%      $$\frac{\delta \sqrt{\log k}}{\sqrt{\log k} + 4 \delta \sqrt{d}}  \geq \frac{1}{8 \sqrt{cd}} \geq \frac{\varepsilon \sqrt{c}}{8 \delta^2}. $$
%      Rest of the proof of~\Cref{cl:ul} follows without any further changes (where we replace $c$ by $\sqrt{c}$). 
    
% \end{proof}

\section{$\kolp$ algorithm for Latent Polytopes using Singular Value Decomposition}
\label{sec:lkp}

~\Cref{lower-bound-0} showed that a solution to $\kolp$ problem requires the error parameter $\varepsilon$ to be $O^*(1/\sqrt{d})$. \Cref{upper-bound-1} and~\Cref{upper-bound-2} give algorithms achieving this bound. However, for many polytopes with $k$ vertices, we can solve the $\kolp$ problem with $\varepsilon$ being $O^*(1/\sqrt{k})$. However, if $k<d$, this error is too high. To tackle this, we find a good approximation to the subspace spanned by the vertices of $K$, 
then we project to this subspace and use the result in~\Cref{upper-bound-1}. 
One such example is the ``Latent $k-$ Polytope'' (abbreviated $\lkp$) problem which we now describe.

The $\lkp$ problem has been studied in~\cite{BhattacharyyaK20}. Certain assumptions were made on the model, namely, the hidden polytope $K$ as well as on the (hidden) process for generating observed data from latent points in $K$. These assumptions are (a) shown to hold in several important Latent Variable models and (b) are sufficient to enable one to get polynomial time learning algorithms.

Here, we formulate assumptions which are similar, but, weaker in one important aspect. Whereas \cite{BhattacharyyaK20} assumed that each vertex of $K$ has a separation from the {\bf affine hull} of the other vertices (thus, in particular, each vertex is affinely independent of other vertices), we assume here that each vertex is separated only from the convex hull of the others. Under this weaker assumption, the algorithm of~\cite{BhattacharyyaK20} does not work. We give a different algorithm which we prove works. It is also simpler to state and carry out and its proof is based on a new general tool we introduce here - the Random Separating Hyperplane theorem (\Cref{thm:sephp}).  

\noindent
{\bf Assumptions on data in the $\lkp$ problem:} 
Let $\Mdot{1}, \ldots, \Mdot{k}$ denote the vertices of $K$ and $\bM$ be the $d \times k$ matrix with columns representing the vertices of $K$.  
We assume there are latent (hidden) points $P_{\cdot ,j}, j=1,2,\ldots ,n$ in $K$ and observed data points $A_{\cdot,j}, j=1,2,\ldots ,n$ are generated (not necessarily under any stochastic assumptions) by adding {\it displacements} $A_{\cdot,j}-P_{\cdot,j}$ respectively to $P_{\cdot,j}$. Clearly if the displacements are arbitrary, it is not possible to learn $K$ given only the observed data. So we need some bound on the displacements. 

Secondly, if all (or almost all) latent points lie in (or close to) the convex hull of a subset of $k-1$ or fewer vertices of $K$, the missing vertex cannot be learnt. To avoid this, we will assume that there is a certain $w_0$ fraction of latent points close to every vertex of $K$. 

Let \footnote{By the standard definition of spectral norm, it is easy to see that $\sigma_0^2$ is the maximum mean squared displacement in any direction.}
\begin{equation}
\label{eq:p}
\sigma_0 :=\frac{||\bP-\bA||}{\sqrt n}.
\end{equation}
We now show that the {\bf $\kolp$-Algorithm} mentioned in~\Cref{algorithmic-rsh}
has the desired properties: 
\begin{theorem}\label{LkP-Main-Thm}
Suppose $K$ is a latent poytope with $k$ vertices $M_{\cdot,1},M_{\cdot,2},\ldots ,M_{\cdot, k}$ and $\bP,\bA$ are latent points (all in $K$) and observed data respectively. Assume 
\begin{equation}\label{eq:cell}
\mbox{For all }\ell\in[k]\; ,\; C_\ell :=\{j:|P_{\cdot,j}-M_{\cdot,\ell}|\leq\frac{\sigma_0}{\sqrt{w_0}}\} \mbox{ satisfies    }
|C_\ell|\geq w_0n.
\end{equation}
Suppose $(\sqrt{\log k}/\sqrt{c_0k})\leq \delta\leq 1$ and $c_0$ is a large constant satisfying
\begin{equation}\label{1000}
\sigma_0\leq \frac{\delta^2\Delta(K)}{100c_0}\frac{\sqrt{w_0}}{\sqrt k}.
\end{equation}
Let $V$ be the $k$-dimensional SVD subspace of $\bf A$, and $\hatK$ denote the projection of $K$ on $V$. 
\begin{itemize}
    \item There is an $ \opto_{\frac{10\sigma_0}{\sqrt{w_0}\Delta}}(\hatK)$ oracle $\cal O$. 
    \item The algorithm {\bf $\kolp$ Algorithm} in~\Cref{algorithmic-rsh} outputs a set $Q$ of $k$ points such that the following condition is satisfied w.h.p.:  for every vertex $v$ of $K$, there is a point $v'$ in $Q$ with $|v-v'| \leq \delta \Delta(K)/5$. 
\end{itemize}
\end{theorem}

% {\color{red} The theorem needs to have a proof. Should be short based on existing claims and Lemmas. Pl put in.   ravi}
In the Theorem above, \eqref{1000}  implies an upper bound on $\sigma_0$ of $\Delta(K)\sqrt{w_0}/(c_0\sqrt k)$ and so the  oracle $\cal O$ used by the Theorem lies in $\opt_\varepsilon(\hatK)$ for  $\varepsilon\leq 1/(c_0\sqrt k)$. Thus, we get around the lower bound result~\Cref{lower-bound-0} by working in the $k$-dimensional SVD subspace of $\bA$. In rest of this section, we prove~\Cref{LkP-Main-Thm}. We begin by giving properties of the SVD subspace, and then show that an approximate optimization oracle exists in this subspace. Finally, we apply~\Cref{upper-bound-2kolp}.



%%%%%%%%%%%

% \subsection{Assumptions on Model $\bM$ and data $\bA$}
% The $\lkp$ problem has been studied in~\cite{BhattacharyyaK20}. Certain assumptions were made on the model, namely, the hidden polytope $K$ as well as on the (hidden) process for generating observed data from latent points in $K$. These assumptions are (a) shown to hold in several important Latent Variable models and (b) are sufficient to enable one to get polynomial time learning algorithms.

% Here, we formulate assumptions which are similar, but, weaker in one important aspect. Whereas~\cite{BhattacharyyaK20} assumed that each vertex of $K$ has a separation from the {\bf affine hull} of the other vertices (thus, in particular, each vertex is affinely independent of other vertices), we assume here that each vertex is separated only from the convex hull of the others. Under this weaker assumption, the algorithm of~\cite{BhattacharyyaK20} does not work. We give a different algorithm which we prove works. It is also simpler to state and carry out and its proof is based on a new general tool we introduce here - the Random Separating Hyperplane theorem (\Cref{thm:sephp}).  

% \noindent
% {\bf Assumptions on the data:} 
% Let $\Mdot{1}, \ldots, \Mdot{k}$ denote the vertices of $K$ and $\bM$ be the $d \times k$ matrix with columns representing the vertices of $K$.  
% We assume there are latent (hidden) points $P_{\cdot ,j}, j=1,2,\ldots ,n$ in $K$ and observed data points $A_{\cdot,j}, j=1,2,\ldots ,n$ are generated (not necessarily under any stochastic assumptions) by adding {\it displacements} $A_{\cdot,j}-P_{\cdot,j}$ respectively to $P_{\cdot,j}$. Clearly if the displacements are arbitrary, it is not possible to learn $K$ given only the observed data. So we need some bound on the displacements. This is done by upper bounding $\sigma_0$ defined below.

% Secondly, if all (or almost all) latent points lie in (or close to) the convex hull of a subset of $k-1$ or fewer vertices of $K$, the missing vertex cannot be learnt. To avoid this, we will assume that there is a certain $w_0$ fraction of latent points close to every vertex of $K$. 

% Let \footnote{By the standard definition of spectral norm, it is easy to see that $\sigma_0^2$ is the maximum mean squared displacement in any direction.}
% \begin{equation}
% \label{eq:p}
% \sigma_0 :=\frac{||\bP-\bA||}{\sqrt n}.
% \end{equation}
% We assume there is a real number $w_0\in (0,1]$ such that (with $c_0$ as in (\ref{c0})):
% {\color{red} Or just say, $c_0$ large enough constant???}
% \begin{equation}
% \label{eq:separation}
%  \sigma_0\leq \frac{\sqrt{w_0}\Delta (K)}{9c_0\sqrt k}.  
% \end{equation}

% \noindent
% We also assume
% \begin{equation}\label{eq:cell}
% \mbox{For all }\ell\in[k]\; ,\; C_\ell :=\{j:|P_{\cdot,j}-M_{\cdot,\ell}|\leq\frac{\sigma_0}{\sqrt{w_0}}\} \mbox{ satisfies    }
% |C_\ell|\geq w_0n.
% \end{equation}

% \noindent
% {\bf Assumption on the model}: we assume that the vertices of $K$ are $\delta$-well separated, where, $\delta $ satisfies: (cf (\ref{eq:separation}):
% \begin{equation}\label{delta}
% 1\geq \delta \geq \mbox{Max}\left( \sqrt{\frac{\ln k}{c_0k}}\; ,\; 
% 3 \left(\frac{k}{w_0}\right)^{1/4}\sqrt{c_0}\left(\frac{\sigma_0}{\Delta(K)}\right)^{1/2}\right).
% \end{equation}
% with $c_0$ a large enough positive constant.

% %%%%%%%%%%

% % \noindent
% % \paragraph{\bf Standard Assumptions on $K$:}
% % We now state the conditions satisfied by such a polytope $K$ in $\Re^d$. 
% % %As before, we can assume wlog that $\Delta(K)$, the diameter of $K$, is exactly 1. 
% % Let $\Mdot{1}, \ldots, \Mdot{k}$ denote the vertices of $K$ and $\bM$ be the $d \times k$ matrix with columns representing the vertices of $K$.  We are also given points $\Adot{1}, \ldots, \Adot{n}$ in $\Re^d$, and let $\bA$ be the $d \times n$ matrix representing these points (note that these points are not in $K$, but are supposed to be {\em perturbations} of points in $K$). 


% % The conditions are as follows:
% % \begin{itemize}
% %     \item[{\bf (A1)}] For each point $\Adot{i}$, there is a point $\Pdot{i} \in K$. Let $\bP$ be the $d \times n$ matrix with columns $\Pdot{1}, \ldots, \Pdot{n}$ (we shall call these ``latent'' points). Let $\sigma_0$ denote the parameter given by
% %     \begin{align}
% %     \label{eq:p}
% %     \frac{||P-A||}{\sqrt{n}} = \frac{\sigma_0}{\Delta(K)}. 
% %     \end{align} This captures the fact that $\bA$ is a perturbation of $\bP$ with bounded spectral norm. 
% %     \item[{\bf (A2)}] There is a parameter $w_0$ such that for every vertex $\Mdot{\ell}$ of $K$, the set of latent points {\em close} to $\Mdot{\ell}$ has size at least $w_0 n,$ i.e., define 
% %     \begin{align}
% %     \label{eq:cell}
% %     C_\ell := \{j : |\Pdot{j} - \Mdot{\ell}| \leq \frac{\sigma_0}{\sqrt{w_0}} \}.
% %     \end{align}
% %     Then for all $\ell\in [k], $ $|C_\ell| \geq w_0 n$. Since the sets $C_\ell$ would turn out to be disjoint, the parameter $w_0$ will be at most $1/k$. 
% %     \item[{\bf (A3)}] The polytope $K$ is $\delta$-well-separated for some parameter $\delta.$ Further, there is a large enough constant $c_0$ such that 
% %     \begin{equation}\label{eq:separation}
% % \delta^2 \geq \max \left( \frac{c_0\sqrt k\sigma_0}{\sqrt{w_0}}\; ,\; \frac{\log k}{c_0 k} \right).
% % \end{equation}
% % \end{itemize}

\subsection{Properties of the SVD subspace of $\bA$}
%For sake of brevity assume that $\Delta(K)=1$. 
Let $V$ denote  the $k$-dimensional SVD subspace of $\bA$. We shall show that the projection of $K$ on $V$ has an   $\opto_\varepsilon(K)$ oracle for a suitable value of $\varepsilon$. Let $\hMdot{\ell}$ denote the projection of $\Mdot{\ell}$ on the SVD subspace $V$. Define $\hAdot{j}$ similarly, and let $\widehat{\bA}$ be the $d \times n$ matrix whose columns are given by $\hAdot{j}$. We now show that $\hMdot{\ell}$ and  $\Mdot{\ell}$ are close to each other. The following notation will turn out to be very useful in subsequent discussion: for a $d \times n$ matrix ${\bf B}$ and a subset $S$ of $[n]$, define 
$$ B_{\cdot, S} := \frac{1}{|S|} \sum_{j \in S} B_{\cdot, j}.$$

\begin{claim}
\label{cl:sum}
%Assume $\cE$ occurs. 
Let $S \subseteq [n]$ be a subset of indices. Then, $$|A_{\cdot, S} - P_{\cdot, S}| \leq \frac{\sigma_0 \sqrt n}{\sqrt{|S|}}. $$
\end{claim}
\begin{proof}
Define a unit vector $v \in \Re^d$ as follows: if $j \in S$, we define $v_j = \frac{1}{\sqrt{|S|}}$, 0 otherwise. Now, observe that 
$$|A_{\cdot, S} - P_{\cdot, S}| = \frac{1}{\sqrt{|S|}} \cdot |(\bA - \bP) \cdot v|. $$
Inequality~\eqref{eq:p} now implies that the RHS above is at most 
$\frac{\sigma_0 \sqrt{n}}{\sqrt{|S|}}.$
\end{proof}


\begin{lemma}
\label{lem:close}
%Assume that the event $\cE'$ happens. 
For all $\ell\in [k]$, 
$|\Mdot{\ell}-\hMdot{\ell}|\leq 5 \sigma_0/\sqrt{w_0} \leq \frac{\delta^2 \Delta(K)}{c_0}$.
\end{lemma}

\begin{proof}
We have $||\widehat{\bA}-\bA||\leq ||\bA-\bP||$ since, $\widehat{\bA}$ is the best rank $k$ approximation to $\bA$ in terms of the spectral norm and since each column of $\bP$ being a convex combination of the columns $\{ \Mdot{\cdot,\ell},\ell\in [k]\}$,
$\bP$ has rank at most $k$. We also have $||\widehat{\bA}-\widehat{\bP}||\leq ||\bA-\bP||$ since projections cannot increase length. Using these inequalities and the triangle inequality, we get: 
$$||\bP-\widehat{\bP}||\leq ||\bP-\bA||+||\bA-\widehat{\bA}||+||\widehat{\bA}-\widehat{\bP}||\leq 3 \cdot ||\bA-\bP|| \stackrel{\eqref{eq:p}}{\leq} 3\sigma_0\sqrt n.$$ 

Let $C_\ell$ be as in~\eqref{eq:cell}.  
 Now, define $w$ to be the unit vector with $w_j=1/\sqrt{|C_\ell|}, \forall j\in C_\ell$, and $w_j=0, \forall j\notin C_\ell$. We see that 
\begin{align*}
& |\Mdot{\ell}-\hMdot{\ell}|\leq \frac{|(\bP-\widehat{\bP})w|}{|C_\ell|}+\frac{2\sigma_0}{\sqrt{w_0}} 
\leq \frac{||\bP -\widehat {\bP }||}{\sqrt{|C_\ell |}} +\frac{2\sigma_0}{\sqrt{w_0}},
\end{align*}
 which proves the Lemma (using $|C_\ell| \geq w_0 n$). The second inequality in the Lemma follows from~\eqref{1000}.
\end{proof}

Let $\hatK$ be the projection of $K$ on $V$. We now show that $\hatK$ is also $\delta'$-well-separated for a value $\delta'$ close to $\delta$. 
\begin{lemma}
    \label{lem:hatK}
    The vertices of $\hatK$ are given by  $\{\hMdot{\ell}: \ell \in [k]\}$. Further, $ \hatK$ is $\delta'$-well-separated, where $\delta'= \delta \left( 1 - \frac{1}{100} \right)$. 
\end{lemma}
\begin{proof}
    It is clear that the vertex set of $\hatK$ is s subset of $S := \{\hMdot{\ell}: \ell \in [k]\}$. Thus, we need to show that none of the points in $S$ can be written as a convex combination of rest of the points in $S$. This will follow from the fact that $\hatK$ is $\delta'$-well-separated, and so it suffices to prove this statement. Consider a point $\hMdot{\ell} \in S$ and a point $p:= \sum_{\ell' \in [k], \ell' \neq \ell} \lambda_{\ell'} \hMdot{\ell'},$ which is in $\CH(S \setminus\{\hMdot{\ell}\}).$ Then 
    \begin{align*}
        |p-\hMdot{\ell}| & = \left| \sum_{\ell' \neq \ell} \lambda_{\ell'} (\hMdot{\ell'} - \hMdot{\ell} ) \right| \\
        & \geq \left| \sum_{\ell' \neq \ell} \lambda_{\ell'} (\Mdot{\ell'} - \Mdot{\ell} ) \right| - \sum_{\ell' \neq \ell} \lambda_{\ell'} |\Mdot{\ell'}-\hMdot{\ell'}| - 
       \sum_{\ell' \neq \ell} \lambda_{\ell'} |\Mdot{\ell} - \hMdot{\ell}| \\
        & = \left| \sum_{\ell' \neq \ell} \lambda_{\ell'} \Mdot{\ell'} - \Mdot{\ell} \right| - \frac{2\delta^2 \Delta(K)}{c_0} \\
        & \geq (\delta - \frac{2\delta^2}{c_0}) \Delta(K) \geq \delta' \Delta(K)
    \end{align*}
    where the second last line uses~\Cref{lem:close} and last line follows from  the fact that $K$ is $\delta$-well-separated. Note that $\Delta(\hatK) \leq \Delta(K).$ Thus it follows that $\hatK$ is is $\delta'$-well-separated.
\end{proof}

% {\color{red} Below in several places, $\varepsilon$ has dimensions of length. It really should be dimensionless. Think, we just have to divide the values set for $\varepsilon$ by $\Delta(K)$ for all occurences of $\varepsilon$ till the end of the section.}


We would now like to use~\Cref{upper-bound-2} on $\hatK$ with $\varepsilon = \frac{\sigma_0}{\sqrt{w_0}\Delta(K)}$ and $\delta'=\delta$. Indeed the condition~\eqref{1000} along with~\Cref{lem:hatK} show that the conditions of~\Cref{upper-bound-2} hold provided we are able to exhibit an efficient $\opto_\varepsilon(\hatK)$ oracle. 

\subsection{Construction of the $\opto_\varepsilon(\hatK)$ oracle}

We now describe the construction for the   $\opto_\varepsilon(\hatK)$ oracle in the subspace $V$, where $\varepsilon = \frac{10 \sigma_0}{\Delta(K) \sqrt{w_0}}$. Let $u$ be a unit vector in $V$. The procedure (referred as {\bf Subset Smoothing Algorithm} in~\Cref{algorithmic-rsh}) is given in~\Cref{algo:center}  -- we project the points in $\hatA$ on $u$ and consider the  $w_0 n$ with the highest projection along $u$. Finally, we output the average of these points. 



\begin{algorithm}[H]
  \caption{Orcale in $\opto_\varepsilon(\hatK)$}
  \label{algo:center}
    {\bf Input:} unit vector $u \in V$. \;
       Let $R(u)$ be the index set of the columns $\hAdot{j}$ with the $w_0 n$ highest values of $u \cdot \hAdot{j}$. \;
      {\bf Output:} $\hAdot{R(u)}$ \;
\end{algorithm}


We now show that this algorithm has the desired properties. Let $R(u)$ be the index set as in~\Cref{algo:center}. The following is an easy consequence of~\Cref{cl:sum}. 
\begin{claim}
\label{cl:eps1}
    $\dst(\hAdot{R(u)}, \hatK) \leq \varepsilon\Delta(K) . $
\end{claim}
\begin{proof}
    \Cref{cl:sum} shows that $|\Adot{R(u)} - \Pdot{R(u)}| \leq \varepsilon\Delta(K) .$ Therefore, 
    $|\hAdot{R(u)} - \hPdot{R(u)}| \leq \varepsilon\Delta(K) $. The desired result now follows because
    $\hPdot{R(u)} \in \hatK.$
\end{proof}


Let $\ell \in [k]$ be the index for which $\hMdot{\ell} \cdot u$ is maximized. 
\begin{claim}
    \label{cl:eps2}
    $\hAdot{R(u)} \cdot u \geq \hMdot{\ell} \cdot u - \varepsilon\Delta(K) .$
\end{claim}
\begin{proof}
    Let $C_\ell$ be the index set specified by~\eqref{eq:cell}. It suffices to show that 
    $$ \hAdot{C_\ell} \cdot u \geq \hMdot{\ell} \cdot u - \varepsilon\Delta(K) .$$ 
    Now, 
    \begin{align*}
    \hMdot{\ell} \cdot u - \hAdot{C_\ell} \cdot u & \leq |\hMdot{\ell} - \hAdot{C_\ell}| \\
    & \leq |\hMdot{\ell} - \hPdot{C_\ell}| + |\hPdot{C_\ell} - \hAdot{C_\ell}| \\
    & \leq |\Mdot{\ell} - \Pdot{C_\ell}| + |\Pdot{C_\ell} - \Adot{C_\ell}| \\
    & \leq \varepsilon \Delta(K),
    \end{align*}
where the last inequality follows from~\eqref{eq:cell} and~\Cref{cl:sum}. 
\end{proof}

The above two results show that~\Cref{algo:center} yields an  $\opto_\varepsilon(\hatK)$ oracle.  We would now like to apply~\Cref{upper-bound-2kolp} to the polytope $\hatK$ with parameters $\varepsilon$ and $\delta' = \delta(1-1/100).$~\Cref{lem:hatK} shows that $\hatK$ is $\delta'$-well-separated. We now need to check that the parameters $\varepsilon$ and $\delta'$ satisfy the following conditions needed in the statement of~\Cref{upper-bound-2kolp} (recall that $\hatK$ is a polytope in $\Re^k$): (i)  $\delta^2 \geq c \varepsilon \sqrt{k}$, and (ii) $\delta \geq \frac{\sqrt{\log k}}{\sqrt{c} k}.$
The second condition is already an assumption in the statement of~\Cref{LkP-Main-Thm}, and the first condition follows from~\eqref{1000} and the fact that $\varepsilon = \frac{10 \sigma_0}{\Delta(K) \sqrt{w_0}}$. Applying~\Cref{upper-bound-2kolp}, we get a set of $k$ points $Q$, such that for each vertex $v'' $ of $\hatK$, there is a point $v' \in Q$ with $|v''-v'| \leq \delta \Delta(K')/10$. Applying~\Cref{lem:close} proves~\Cref{LkP-Main-Thm}. 

